% ============================================================
%  sections/spill_response.tex
%  Section 3 — Spill Response and Emergency Procedures
%  Includes the TikZ spill-response flowchart
% ============================================================

\section{Spill Response and Emergency Procedures}
\label{sec:spill}

Biological spills at BSL-2 must be managed calmly and systematically. The response
differs depending on whether aerosolisation is likely (e.g., centrifuge failure,
pressurised container breach) or unlikely (e.g., overturned culture flask on bench).
The following flowchart provides decision logic for both scenarios.

% -----------------------------------------------------------
\subsection{Spill Response Flowchart}
% -----------------------------------------------------------

\vspace{8pt}
\begin{figure}[H]
\centering
\begin{tikzpicture}[
  node distance = 0.55cm and 1.0cm,
  startstop/.style = {
    rectangle, rounded corners=5pt,
    minimum width=4.0cm, minimum height=0.82cm,
    text centered, draw=BSLdark, fill=BSLdark,
    font=\sffamily\bfseries\footnotesize\color{white}
  },
  process/.style = {
    rectangle,
    minimum width=4.4cm, minimum height=0.82cm,
    text width=4.2cm,
    text centered, draw=BSLmid, fill=BSLtabalt,
    font=\sffamily\footnotesize\color{BSLdark}
  },
  decision/.style = {
    diamond, aspect=2.6,
    minimum width=4.8cm, minimum height=0.9cm,
    text width=3.8cm,
    text centered, draw=BSLamber, fill=BSLwarn,
    font=\sffamily\bfseries\footnotesize\color{BSLdark}
  },
  danger/.style = {
    rectangle,
    minimum width=4.4cm, minimum height=0.82cm,
    text width=4.2cm,
    text centered, draw=BSLdangerborder, fill=BSLdanger,
    font=\sffamily\footnotesize\color{BSLdangerborder}
  },
  arrow/.style = {->, thick, color=BSLmid},
  label/.style = {font=\sffamily\tiny\color{BSLmid}}
]

% ---- nodes ----
\node (spill)    [startstop]                          {SPILL DETECTED};
\node (d1)       [decision, below=of spill]           {Aerosol likely?\\\scriptsize(centrifuge, pressurised vessel)};
\node (evac)     [danger,   below left=1.2cm and 0.4cm of d1]  {Evacuate room.\newline Alert all personnel.\newline Close door.};
\node (alert)    [process,  below=of evac]            {Call Biosafety Emergency:\newline \textbf{+1 (617) 555-0194}};
\node (wait)     [process,  below=of alert]           {Wait 30 minutes outside for aerosol to settle.};
\node (ppe1)     [process,  below right=1.2cm and 0.4cm of d1] {Don appropriate PPE:\newline coat, double gloves, face shield.};
\node (cover)    [process,  below=of ppe1]            {Cover spill with paper towels.\newline Do not spread material.};
\node (decon)    [process,  below=of cover]           {Apply 10\% bleach around then over spill. Contact 20 min.};
\node (rejoin)   [process,  below=of wait]            {Re-enter with full PPE.\newline Confirm aerosol settled.};
\node (d2)       [decision, below=4.5cm of d1]        {Spill contained \&\newline risk neutralised?};
\node (collect)  [process,  below=of d2]              {Collect material with absorbents.\newline Seal in biohazard bag.};
\node (bso)      [danger,   right=1.4cm of d2]        {Contact BSO immediately.\newline Do NOT proceed alone.};
\node (decon2)   [process,  below=of collect]         {Decontaminate area with fresh 10\% bleach. Clean \& dispose.};
\node (report)   [process,  below=of decon2]          {Complete Incident Report (Form MRI-INC-01) within 24~h.};
\node (done)     [startstop,below=of report]          {INCIDENT CLOSED};

% ---- arrows ----
\draw[arrow] (spill)  -- (d1);
\draw[arrow] (d1)     -| node[label, above, pos=0.25]{Yes} (evac);
\draw[arrow] (d1)     -| node[label, above, pos=0.25]{No}  (ppe1);
\draw[arrow] (evac)   -- (alert);
\draw[arrow] (alert)  -- (wait);
\draw[arrow] (wait)   -- (rejoin);
\draw[arrow] (ppe1)   -- (cover);
\draw[arrow] (cover)  -- (decon);
\draw[arrow] (rejoin) |- (d2);
\draw[arrow] (decon)  |- (d2);
\draw[arrow] (d2)     -- node[label, right]{Yes} (collect);
\draw[arrow] (d2)     -- node[label, above]{No}  (bso);
\draw[arrow] (collect) -- (decon2);
\draw[arrow] (decon2) -- (report);
\draw[arrow] (report) -- (done);

\end{tikzpicture}
\caption{BSL-2 Biological Spill Response Decision Flowchart}
\label{fig:spill_flowchart}
\end{figure}

\vspace{4pt}

% -----------------------------------------------------------
\subsection{Step-by-Step Spill Response Procedures}
% -----------------------------------------------------------

\subsubsection{Minor Bench Spill (No Aerosol Risk)}

\begin{procedurebox}[Minor Bench Spill Protocol]
\begin{enumerate}[leftmargin=1.6em, itemsep=3pt,
                  label=\textbf{\color{BSLmid}\arabic*.}]
  \item \textbf{Alert:} Notify others in the room. If the spill involves a known
        pathogen, verbally announce ``biological spill'' so colleagues can take
        appropriate precautions.
  \item \textbf{PPE:} Don double gloves, safety glasses, and lab coat if not
        already wearing them. Add a face shield if splashing is possible.
  \item \textbf{Contain:} Place dry paper towels over the spill without pressing
        down. Work from the outer edges toward the centre.
  \item \textbf{Disinfect:} Carefully pour 10\% bleach around the perimeter of
        the towels, then over the top. Allow \textbf{20 minutes} contact time.
  \item \textbf{Collect:} Using forceps or a dust pan (never bare hands), collect
        the soaked towels and place in a biohazard bag.
  \item \textbf{Re-clean:} Wipe the area with fresh 10\% bleach, followed by 70\%
        ethanol to remove bleach residue from equipment surfaces.
  \item \textbf{Report:} Document the spill using Form~MRI-INC-01. Minor spills
        not involving personnel exposure may be reported to the PI within 24~h.
\end{enumerate}
\end{procedurebox}

\subsubsection{Major Spill or Centrifuge Accident (Aerosol Risk)}

\begin{dangerbox}[Aerosol-Generating Spill — Immediate Actions]
\small
\textbf{Do not attempt cleanup until the room has been evacuated for at least
30 minutes} to allow aerosols to settle and/or be removed by the ventilation system.
Failure to observe this wait time may result in inhalation exposure.
\end{dangerbox}

\begin{procedurebox}[Aerosol Spill Protocol]
\begin{enumerate}[leftmargin=1.6em, itemsep=3pt,
                  label=\textbf{\color{BSLmid}\arabic*.}]
  \item \textbf{Stop work. Evacuate immediately.} Hold breath and exit. Post
        ``Do Not Enter -- Biological Spill'' sign on the door.
  \item \textbf{Remove contaminated clothing} at the door threshold; bag and seal
        immediately. Wash exposed skin and eyes for 15~minutes at the nearest
        eyewash / safety shower.
  \item \textbf{Call the Biosafety Emergency Line} (+1 617-555-0194) and your
        Principal Investigator within 15~minutes.
  \item \textbf{Wait 30 minutes} before re-entry (aerosol settling time).
  \item \textbf{Re-enter with full PPE}: double gloves, N95 respirator, face
        shield, solid-front gown, boot covers.
  \item \textbf{Proceed as Minor Spill steps 3--7} for cleanup.
\end{enumerate}
\end{procedurebox}

% -----------------------------------------------------------
\subsection{Medical Emergencies and Exposure Incidents}
% -----------------------------------------------------------

\begin{infobox}[Exposure Incident Response]
\small
\begin{tabular}{@{}L{3.2cm} L{10cm}@{}}
\textbf{Needlestick / cut:} & Remove gloves; wash thoroughly with soap and water for
                               5--10~min. Do not squeeze or suction the wound. Report
                               immediately to Employee Health (Building~1, Room~22). \\[4pt]
\textbf{Eye splash:}         & Irrigate with eyewash for minimum 15~min. Report to
                               Employee Health and Biosafety Office immediately. \\[4pt]
\textbf{Ingestion:}          & Do not induce vomiting. Contact Poison Control
                               (1-800-222-1222) and Employee Health. \\[4pt]
\textbf{All exposures:}      & Must be reported within 4~hours using Form~MRI-INC-01.
                               Post-exposure prophylaxis (PEP) availability is agent-specific;
                               Employee Health will advise. \\
\end{tabular}
\end{infobox}

\begin{warningbox}[Never delay seeking medical attention to complete paperwork.]
Form MRI-INC-01 can be completed after receiving initial medical evaluation.
The priority is your health and safety.
\end{warningbox}
